\documentclass[11pt,a4paper]{article}
\usepackage[T1]{fontenc}
\usepackage[utf8]{inputenc}
\usepackage{amsmath,amssymb,xcolor}
\usepackage[margin=25mm]{geometry}
\usepackage[colorlinks=true,urlcolor=blue]{hyperref}
\definecolor{nexorablue}{HTML}{2459B5}
\setlength{\parindent}{0pt}
\setlength{\parskip}{7pt}
\setlength{\emergencystretch}{2em}
\renewcommand{\familydefault}{\sfdefault}
\begin{document}

{\small\bfseries\color{nexorablue} NEXORA CPGE}\par
{\small MPSI / NEX-RGT-05 / Corrigé}\par
{\LARGE\bfseries Régler une décharge au régime critique}\par
\vspace{0.25em}\hrule\vspace{0.6em}

\section*{Valeur critique}
On trouve $\omega_0=500\,\mathrm{s^{-1}}$, puis $R_c=2L\omega_0=40\,\mathrm\Omega$.

\section*{Conditions initiales}
Posons $q(t)=(A+Bt)e^{-500t}$. On a $A=q(0)$ et $q'(0)=B-500A=0$. Donc $q(t)=q(0)(1+500t)e^{-500t}$. La charge tend vers zéro sans changer de signe.


\end{document}
